%%%%%%%%%%%%%%%%%%%%%%%%%%%%% Define Article %%%%%%%%%%%%%%%%%%%%%%%%%%%%%%%%%%
\documentclass{article}
%%%%%%%%%%%%%%%%%%%%%%%%%%%%%%%%%%%%%%%%%%%%%%%%%%%%%%%%%%%%%%%%%%%%%%%%%%%%%%%

%%%%%%%%%%%%%%%%%%%%%%%%%%%%% Using Packages %%%%%%%%%%%%%%%%%%%%%%%%%%%%%%%%%%
\usepackage{listings, xcolor}
\usepackage{parskip}
\usepackage{hyperref}
\usepackage{amsmath}
\usepackage{graphicx}
\usepackage[a4paper, margin = 1.5cm]{geometry}
\usepackage{float}
\usepackage{amsmath}
\usepackage{amssymb}
\usepackage{parskip}

%%%%%%%%%%%%%%%%%%%%%%%%%%%%%%% Title & Author %%%%%%%%%%%%%%%%%%%%%%%%%%%%%%%%
\title{Introduzione Basi di Dati}
\author{Trentin Tommaso}
%%%%%%%%%%%%%%%%%%%%%%%%%%%%%%%%%%%%%%%%%%%%%%%%%%%%%%%%%%%%%%%%%%%%%%%%%%%%%%%

\begin{document}
    \maketitle
    \tableofcontents
    \newpage

    \section{Modelli di processi SW}
      \subsection{Processo vs Modello di processo}
        Non esiste modello univesale $\rightarrow$ ciascuna organizzazione usa proprio processo (possibilità di usare processi diversi per prodotti diversi).\\
        Ogni processo comprende attività, sotto-attività, prodotti e persone con responsabilità.\\
        \textbf{Modello di processo SW:} descrizione semplificata/astratta del processo SW, ci si riferisce a questi modelli con diversi alias (modelli di processo sw / ciclo di vita del sw / paradigmi di processo).\\
        Descrizioni astratte di alto livello di processi sw che possono essere utilizzate per spiegare diversi approcci allo sviluppo sw.\\
      \subsection{Modello si processo SW: Punti di vista}
        Ogni processo ha diverse descrizioni possibili:
        \begin{itemize}
          \item Architetturale $\rightarrow$ descrive sequenza di attività senza fornire dettgli su specifiche attività;
          \item Data-flow $\rightarrow$ evidenzia trasformazioni dei dati operate da attività del processo;
          \item Role/action $\rightarrow$ ruoli delle persone coinvolte nel processo e relative responsabilità;
        \end{itemize}
      \subsection{Descrizione del processo SW}
        Descrizione del processo include:
        \begin{itemize}
          \item \textbf{Attività} da svolgere e loro ordine;
          \item \textbf{Prodotti} risultanti da ciascuna attività;
          \item \textbf{Ruoli} (persone coinvolte e loro responsabilità)
          \item \textbf{Pre/Post condizioni} che devono verificarsi prima/dopo un'attività del processo sia svolta. Stabiliscono criteri di transizione per progredire da uno stadio del processo al successivo.
        \end{itemize}
\end{document}
